
\documentclass{article}
\usepackage{colortbl}
\usepackage{makecell}
\usepackage{multirow}
\usepackage{supertabular}

\begin{document}

\newcounter{utterance}

\centering \large Interaction Transcript for game `wordle', experiment `high\_frequency\_words\_no\_clue\_no\_critic', episode 5 with gpt{-}5.4{-}azure.
\vspace{24pt}

{ \footnotesize  \setcounter{utterance}{1}
\setlength{\tabcolsep}{0pt}
\begin{supertabular}{c@{$\;$}|p{.15\linewidth}@{}p{.15\linewidth}p{.15\linewidth}p{.15\linewidth}p{.15\linewidth}p{.15\linewidth}}
    \# & \multicolumn{2}{c}{Player} && \multicolumn{2}{c}{Game Master} \\
    \hline

    \theutterance \stepcounter{utterance}  
    & & & \multicolumn{4}{p{0.6\linewidth}}{
        \cellcolor[rgb]{0.9,0.9,0.9}{
            \makecell[{{p{\linewidth}}}]{
                \texttt{\tiny{[P1$\langle$GM]}}
                \texttt{Je bent een taaltovenaar die graag woorden raadt met behulp van de gegeven regels.} \\
\\ 
\texttt{Welkom bij Wordle! Je hebt zes pogingen om het doelwoord te raden, een geldig woord van vijf kleine letters. Gebruik de tags "uitleg:" en "gok:" om voor elke gok een beknopte uitleg te geven.} \\
\\ 
\texttt{Bijvoorbeeld, als je gok "staat" is, moet je antwoord zijn} \\
\texttt{uitleg: dit is een veelvoorkomend woord van vijf letters, en ik begin mijn gok met dit woord.} \\
\texttt{gok: staat} \\
\\ 
\texttt{Na elke gok wordt je antwoord gevalideerd en ontvang je feedback die aangeeft welke letters correct zijn (groen), welke letters correct zijn maar op de verkeerde positie staan (geel), en welke letters onjuist zijn (rood). Deze feedback kan nuttig zijn om te bepalen welke letters je in je volgende gok moet opnemen of juist moet uitsluiten.} \\
\\ 
\texttt{Bijvoorbeeld, de feedback voor "staat" kan zijn:} \\
\texttt{gok\_feedback: s<geel> t<geel> a<groen> a<geel> t<rood>} \\
\\ 
\texttt{De uitleg moet details bevatten over hoe de gok\_feedback wordt gebruikt om tot een nieuwe gok te komen.} \\
\\ 
\texttt{Genereer geen andere tekst.} \\
\\ 
\texttt{Laten we beginnen met je eerste gok.} \\
            }
        }
    }
     \\ \\

    \theutterance \stepcounter{utterance}  
    & \multicolumn{4}{p{0.6\linewidth}}{
        \cellcolor[rgb]{0.561,0.737,0.561}{
            \makecell[{{p{\linewidth}}}]{
                \texttt{\tiny{[P1$\rangle$GM]}}
                \texttt{uitleg: dit is een veelvoorkomend woord van vijf letters met meerdere veelgebruikte klinkers en medeklinkers, zodat de eerste feedback veel informatie geeft.} \\
\texttt{gok: stoel} \\
            }
        }
    }
    & & \\ \\

    \theutterance \stepcounter{utterance}  
    & & & \multicolumn{2}{p{0.3\linewidth}}{
        \cellcolor[rgb]{0.9,0.9,0.9}{
            \makecell[{{p{\linewidth}}}]{
                \texttt{\tiny{[GM$|$GM]}}
                \texttt{attempts = 1} \\
\texttt{target\_word = sfeer} \\
\texttt{guess = stoel} \\
\texttt{guess\_feedback = s<green> t<red> o<red> e<green> l<red>} \\
            }
        }
    }
    & & \\ \\

    \theutterance \stepcounter{utterance}  
    & & & \multicolumn{4}{p{0.6\linewidth}}{
        \cellcolor[rgb]{0.9,0.9,0.9}{
            \makecell[{{p{\linewidth}}}]{
                \texttt{\tiny{[P1$\langle$GM]}}
                \texttt{gok\_feedback: s<groen> t<rood> o<rood> e<groen> l<rood>} \\
\\ 
\texttt{Geef uw antwoord alleen in dit formaat.} \\
\texttt{uitleg: details} \\
\texttt{gok: woord} \\
            }
        }
    }
     \\ \\

    \theutterance \stepcounter{utterance}  
    & \multicolumn{4}{p{0.6\linewidth}}{
        \cellcolor[rgb]{0.561,0.737,0.561}{
            \makecell[{{p{\linewidth}}}]{
                \texttt{\tiny{[P1$\rangle$GM]}}
                \texttt{uitleg: de feedback laat zien dat s op positie 1 goed is en e op positie 4 ook. De letters t, o en l komen niet in het woord voor, dus ik vermijd die volledig. Ik kies een woord in het patroon s \_ \_ e \_ met nieuwe veelvoorkomende letters om extra informatie te krijgen.} \\
\texttt{gok: spier} \\
            }
        }
    }
    & & \\ \\

    \theutterance \stepcounter{utterance}  
    & & & \multicolumn{2}{p{0.3\linewidth}}{
        \cellcolor[rgb]{0.9,0.9,0.9}{
            \makecell[{{p{\linewidth}}}]{
                \texttt{\tiny{[GM$|$GM]}}
                \texttt{attempts = 2} \\
\texttt{target\_word = sfeer} \\
\texttt{guess = spier} \\
\texttt{guess\_feedback = s<green> p<red> i<red> e<green> r<green>} \\
            }
        }
    }
    & & \\ \\

    \theutterance \stepcounter{utterance}  
    & & & \multicolumn{4}{p{0.6\linewidth}}{
        \cellcolor[rgb]{0.9,0.9,0.9}{
            \makecell[{{p{\linewidth}}}]{
                \texttt{\tiny{[P1$\langle$GM]}}
                \texttt{gok\_feedback: s<groen> p<rood> i<rood> e<groen> r<groen>} \\
\\ 
\texttt{Geef uw antwoord alleen in dit formaat.} \\
\texttt{uitleg: details} \\
\texttt{gok: woord} \\
            }
        }
    }
     \\ \\

    \theutterance \stepcounter{utterance}  
    & \multicolumn{4}{p{0.6\linewidth}}{
        \cellcolor[rgb]{0.561,0.737,0.561}{
            \makecell[{{p{\linewidth}}}]{
                \texttt{\tiny{[P1$\rangle$GM]}}
                \texttt{uitleg: uit de feedback blijkt dat s op positie 1, e op positie 4 en r op positie 5 vaststaan. De letters p en i vallen af, net als eerder t, o en l. Daarom kies ik een woord in het patroon s \_ \_ e r met nieuwe letters die nog niet getest zijn.} \\
\texttt{gok: scher} \\
            }
        }
    }
    & & \\ \\

    \theutterance \stepcounter{utterance}  
    & & & \multicolumn{2}{p{0.3\linewidth}}{
        \cellcolor[rgb]{0.9,0.9,0.9}{
            \makecell[{{p{\linewidth}}}]{
                \texttt{\tiny{[GM$|$GM]}}
                \texttt{Error: The guessed word is not a valid word for this game.} \\
            }
        }
    }
    & & \\ \\

    \theutterance \stepcounter{utterance}  
    & & & \multicolumn{4}{p{0.6\linewidth}}{
        \cellcolor[rgb]{0.9,0.9,0.9}{
            \makecell[{{p{\linewidth}}}]{
                \texttt{\tiny{[P1$\langle$GM]}}
                \texttt{Uw gok is geen geldig woord voor dit spel. Probeer het opnieuw.} \\
\\ 
\texttt{Geef uw antwoord alleen in dit formaat.} \\
\texttt{uitleg: details} \\
\texttt{gok: woord} \\
            }
        }
    }
     \\ \\

    \theutterance \stepcounter{utterance}  
    & \multicolumn{4}{p{0.6\linewidth}}{
        \cellcolor[rgb]{0.561,0.737,0.561}{
            \makecell[{{p{\linewidth}}}]{
                \texttt{\tiny{[P1$\rangle$GM]}}
                \texttt{uitleg: de vorige gok was ongeldig, dus ik probeer een geldig woord dat nog steeds past bij de bekende informatie: s op positie 1, e op positie 4 en r op positie 5. De letters t, o, l, p en i sluit ik uit. Daarom kies ik een geldig woord in het patroon s \_ \_ e r.} \\
\texttt{gok: smeer} \\
            }
        }
    }
    & & \\ \\

    \theutterance \stepcounter{utterance}  
    & & & \multicolumn{2}{p{0.3\linewidth}}{
        \cellcolor[rgb]{0.9,0.9,0.9}{
            \makecell[{{p{\linewidth}}}]{
                \texttt{\tiny{[GM$|$GM]}}
                \texttt{attempts = 3} \\
\texttt{target\_word = sfeer} \\
\texttt{guess = smeer} \\
\texttt{guess\_feedback = s<green> m<red> e<green> e<green> r<green>} \\
            }
        }
    }
    & & \\ \\

    \theutterance \stepcounter{utterance}  
    & & & \multicolumn{4}{p{0.6\linewidth}}{
        \cellcolor[rgb]{0.9,0.9,0.9}{
            \makecell[{{p{\linewidth}}}]{
                \texttt{\tiny{[P1$\langle$GM]}}
                \texttt{gok\_feedback: s<groen> m<rood> e<groen> e<groen> r<groen>} \\
\\ 
\texttt{Geef uw antwoord alleen in dit formaat.} \\
\texttt{uitleg: details} \\
\texttt{gok: woord} \\
            }
        }
    }
     \\ \\

    \theutterance \stepcounter{utterance}  
    & \multicolumn{4}{p{0.6\linewidth}}{
        \cellcolor[rgb]{0.561,0.737,0.561}{
            \makecell[{{p{\linewidth}}}]{
                \texttt{\tiny{[P1$\rangle$GM]}}
                \texttt{uitleg: de feedback bevestigt nu s op positie 1, e op positie 3 en 4, en r op positie 5. De m valt af, samen met de eerder uitgesloten letters t, o, l, p en i. Daarmee blijft het patroon s \_ e e r over, en ik kies een geldig woord dat precies past.} \\
\texttt{gok: sneer} \\
            }
        }
    }
    & & \\ \\

    \theutterance \stepcounter{utterance}  
    & & & \multicolumn{2}{p{0.3\linewidth}}{
        \cellcolor[rgb]{0.9,0.9,0.9}{
            \makecell[{{p{\linewidth}}}]{
                \texttt{\tiny{[GM$|$GM]}}
                \texttt{attempts = 4} \\
\texttt{target\_word = sfeer} \\
\texttt{guess = sneer} \\
\texttt{guess\_feedback = s<green> n<red> e<green> e<green> r<green>} \\
            }
        }
    }
    & & \\ \\

    \theutterance \stepcounter{utterance}  
    & & & \multicolumn{4}{p{0.6\linewidth}}{
        \cellcolor[rgb]{0.9,0.9,0.9}{
            \makecell[{{p{\linewidth}}}]{
                \texttt{\tiny{[P1$\langle$GM]}}
                \texttt{gok\_feedback: s<groen> n<rood> e<groen> e<groen> r<groen>} \\
\\ 
\texttt{Geef uw antwoord alleen in dit formaat.} \\
\texttt{uitleg: details} \\
\texttt{gok: woord} \\
            }
        }
    }
     \\ \\

    \theutterance \stepcounter{utterance}  
    & \multicolumn{4}{p{0.6\linewidth}}{
        \cellcolor[rgb]{0.561,0.737,0.561}{
            \makecell[{{p{\linewidth}}}]{
                \texttt{\tiny{[P1$\rangle$GM]}}
                \texttt{uitleg: de feedback bevestigt s op positie 1, e op positie 3 en 4, en r op positie 5. De letter n valt nu ook af, naast m, p, i, t, o en l. Het woord moet dus het patroon s \_ e e r hebben, en een passend geldig woord is met een nieuwe tweede letter f.} \\
\texttt{gok: sfeer} \\
            }
        }
    }
    & & \\ \\

    \theutterance \stepcounter{utterance}  
    & & & \multicolumn{2}{p{0.3\linewidth}}{
        \cellcolor[rgb]{0.9,0.9,0.9}{
            \makecell[{{p{\linewidth}}}]{
                \texttt{\tiny{[GM$|$GM]}}
                \texttt{attempts = 5} \\
\texttt{target\_word = sfeer} \\
\texttt{guess = sfeer} \\
\texttt{guess\_feedback = s<green> f<green> e<green> e<green> r<green>} \\
            }
        }
    }
    & & \\ \\

    \theutterance \stepcounter{utterance}  
    & & & \multicolumn{2}{p{0.3\linewidth}}{
        \cellcolor[rgb]{0.9,0.9,0.9}{
            \makecell[{{p{\linewidth}}}]{
                \texttt{\tiny{[GM$|$GM]}}
                \texttt{game\_result = WIN} \\
            }
        }
    }
    & & \\ \\

\end{supertabular}
}

\end{document}
